\documentclass[12pt,a4paper]{article}

\usepackage[a4paper,margin=2cm]{geometry}
\usepackage[T2A]{fontenc}
\usepackage[utf8]{inputenc}
\usepackage[russian]{babel}
\usepackage{amsmath,amssymb,mathtools}
\usepackage{array,booktabs,longtable}
\usepackage{xcolor}
\usepackage{hyperref}
\usepackage{tikz}

\hypersetup{
  colorlinks=true,
  linkcolor=blue!55!black,
  urlcolor=blue!55!black,
  pdftitle={Ревью домашней работы},
  pdfauthor={Лёвин Артём Александрович}
}

\setlength{\parindent}{0pt}
\setlength{\parskip}{0.65em}
\renewcommand{\arraystretch}{1.22}

\newcommand{\errlabel}[1]{\textcolor{red}{\textbf{Ошибка:}} \textbf{#1}}
\newcommand{\keyidea}{\textcolor{blue!60!black}{\textbf{Ключевая идея:}}}
\newcommand{\outcome}{\textcolor{teal!60!black}{\textbf{Итог:}}}

\begin{document}

\begin{titlepage}
  \centering
  \vspace*{0.23\textheight}
  {\Huge\bfseries Ревью домашней работы\par}
  \vspace{2cm}
  {\large Автор: Лёвин Артём Александрович\par}
  \vspace{0.7cm}
  {\large 14 сентября 2026 года\par}
  \vfill
  \begin{tikzpicture}[scale=1]
    \draw[blue!45!black, line width=0.8pt]
      (-4,0) .. controls (-2.2,0.7) and (2.2,-0.7) .. (4,0);
  \end{tikzpicture}
  \vspace*{1.2cm}
\end{titlepage}

\newpage
\section*{Сводная таблица}

\small
\begin{longtable}{>{\centering\arraybackslash}p{0.07\linewidth} p{0.17\linewidth} p{0.25\linewidth} p{0.17\linewidth} p{0.17\linewidth}}
\toprule
\textbf{№} & \textbf{Тема} & \textbf{Суть ошибки} & \textbf{Ваш ответ} & \textbf{Правильный ответ} \\
\midrule
\endfirsthead
\toprule
\textbf{№} & \textbf{Тема} & \textbf{Суть ошибки} & \textbf{Ваш ответ} & \textbf{Правильный ответ} \\
\midrule
\endhead
\hyperref[task:3]{3} & Частоты и относительная частота & Перепутано значение признака с частотой; относительная частота найдена неверно. & 1) $4$; 2) $1$ & 1) $1$; 2) $\frac14$ \\
\addlinespace
\hyperref[task:4]{4} & Классическая вероятность & Решение не выполнено: не выделены все исходы и благоприятные исходы. & не указан & $\frac12$ \\
\addlinespace
\hyperref[task:5]{5} & Два последовательных испытания & Решение не выполнено: не выписаны результаты двух бросков монеты. & не указан & $\frac12$ \\
\addlinespace
\hyperref[task:6]{6} & Вероятность противоположного события & Решение не выполнено: не подсчитаны все шары и шары подходящих цветов. & не указан & $\frac7{10}$ \\
\addlinespace
\hyperref[task:8]{8} & Среднее, медиана, мода по таблице & Среднее и медиана указаны неверно, мода не указана. & 1) $4$; 2) $3{,}5$; 3) не указан & 1) $3{,}8$; 2) $4$; 3) $4$ \\
\bottomrule
\end{longtable}
\normalsize

\newpage
\null\vfill
\begin{center}
  {\Huge\bfseries Разбор ошибок}
\end{center}
\vfill\null
\newpage

\phantomsection\label{task:3}
\section*{Задание 3}

\textbf{Текст задания.}
В классе провели опрос о количестве домашних животных. Получена таблица:
\[
\begin{array}{c|cccc}
\text{Количество животных} & 0 & 1 & 2 & 3 \\
\hline
\text{Количество учеников} & 4 & 7 & 5 & 4
\end{array}
\]
Найдите: 1) наиболее часто встречающееся количество домашних животных; 2) относительную частоту результата «2 домашних животных».

\textbf{Ваше решение.}
Промежуточные вычисления не приведены. Записано:
\[
1)\ 4, \qquad 2)\ 1.
\]

\textbf{Ваш ответ.}
$1)\ 4;\quad 2)\ 1.$

\errlabel{неверно прочитана статистическая таблица и неверно найдена относительная частота.}

\textbf{Где именно возникает ошибка.}
В работе нет промежуточных шагов, поэтому нельзя надёжно установить, каким именно рассуждением получены числа $4$ и $1$. Первый достоверно неверный фрагмент --- уже записанный ответ к первому пункту: $4$. Второй достоверно неверный фрагмент --- ответ $1$ к пункту об относительной частоте.

Последний фрагмент, который можно сохранить без исправлений, --- сама исходная таблица: значения признака находятся в верхней строке, а количество учеников с соответствующим числом домашних животных --- в нижней.

\textbf{Почему это неверно.}
В первом пункте требуется не частота, то есть не количество учеников, а \emph{значение признака}: сколько домашних животных встречается чаще всего. Поэтому нужно найти наибольшее число в строке «Количество учеников», а затем посмотреть, какому количеству животных оно соответствует.

Частоты равны $4,7,5,4$. Наибольшая частота --- $7$. Она стоит под значением $1$. Значит, чаще всего встречается именно $1$ домашнее животное. Число $4$, записанное в ответе, находится в строке частот и не является требуемым значением признака.

Во втором пункте требуется \emph{относительная частота}. Она показывает, какую долю от общего числа наблюдений составляет интересующий результат. Для результата «2 домашних животных» благоприятных наблюдений $5$, потому что именно $5$ учеников имеют по два домашних животных. Но делить нужно не на одно из чисел таблицы, а на общее количество опрошенных учеников.

\keyidea сначала определите, что обозначает каждая строка таблицы. Для моды или наиболее частого значения ищите максимальную частоту и возвращайтесь к соответствующему значению признака. Для относительной частоты используйте отношение
\[
\frac{\text{число наблюдений нужного результата}}{\text{общее число наблюдений}}.
\]

\textbf{Исправление от последнего правильного шага.}
Сохраняем данные таблицы и считаем общее число учеников:
\[
4+7+5+4=20.
\]
Наибольшая частота равна $7$, поэтому наиболее часто встречающееся количество домашних животных равно $1$.

Для двух домашних животных частота равна $5$, поэтому относительная частота:
\[
\frac{5}{20}=\frac14.
\]

\textbf{Полное правильное решение.}

1. Рассмотрим частоты: $4,7,5,4$.

2. Самая большая частота --- $7$. Эта частота соответствует значению $1$ в строке «Количество животных». Следовательно, наиболее часто встречается $1$ домашнее животное.

3. Для относительной частоты результата «2 домашних животных» сначала находим объём выборки:
\[
4+7+5+4=20.
\]

4. Два домашних животных имеют $5$ учеников. Поэтому доля таких учеников среди всех опрошенных равна
\[
\frac{5}{20}=\frac14.
\]

\textbf{Проверка.}
Дробь $\frac14$ равна $0{,}25$, то есть $25\%$. Это согласуется с таблицей: из $20$ учеников ровно $5$ имеют по два домашних животных, а $5$ --- это четверть от $20$.

\outcome наиболее часто встречающееся количество домашних животных равно $1$, а относительная частота результата «2 домашних животных» равна $\frac14$.

\textbf{Задача на закрепление.}
В классе опросили учеников о количестве домашних животных. Получили таблицу:
\[
\begin{array}{c|cccc}
\text{Количество животных} & 0 & 1 & 2 & 3 \\
\hline
\text{Количество учеников} & 3 & 9 & 6 & 6
\end{array}
\]
Найдите наиболее часто встречающееся количество домашних животных и относительную частоту результата «2 домашних животных».

\newpage
\phantomsection\label{task:4}
\section*{Задание 4}

\textbf{Текст задания.}
Игральный кубик бросают один раз. Найдите вероятность того, что выпавшее число больше $2$ и при этом отлично от $5$.

\textbf{Ваше решение.}
Записано: «не помню». Математическое решение не приведено.

\textbf{Ваш ответ.}
Не указан.

\errlabel{решение не выполнено.}

\textbf{Где именно возникает ошибка.}
Математически неверного преобразования здесь нет, потому что вычисление не начато. Неполнота возникает до первого необходимого шага: не выписаны возможные результаты броска кубика и не выделены исходы, которые одновременно удовлетворяют двум условиям.

Последний корректный момент --- исходное понимание опыта: кубик бросают один раз. Дальше нужно перейти к пространству равновероятных исходов.

\textbf{Почему это нужно сделать.}
У стандартного игрального кубика шесть равновероятных результатов:
\[
1,2,3,4,5,6.
\]
Фраза «больше $2$ и при этом отлично от $5$» содержит союз «и». Это означает, что подходят только те результаты, которые удовлетворяют \emph{обоим} условиям одновременно.

Сначала числа больше $2$:
\[
3,4,5,6.
\]
Затем исключаем число $5$. Остаются:
\[
3,4,6.
\]
Благоприятных исходов $3$, а всех равновероятных исходов $6$.

\keyidea в классической вероятности при равновероятных исходах сначала выпишите все возможные исходы, затем оставьте только благоприятные и найдите отношение их количества к общему числу исходов.

\textbf{Исправление.}
Всего исходов $6$. Благоприятных исходов $3$: это $3,4,6$. Поэтому
\[
P=\frac{3}{6}=\frac12.
\]

\textbf{Полное правильное решение.}

1. Возможные результаты одного броска кубика:
\[
\{1,2,3,4,5,6\}.
\]

2. Выпавшее число должно быть больше $2$, поэтому временно подходят
\[
\{3,4,5,6\}.
\]

3. Одновременно число должно быть отлично от $5$, поэтому исключаем $5$:
\[
\{3,4,6\}.
\]

4. Всего равновероятных исходов $6$, благоприятных --- $3$. Следовательно,
\[
P=\frac{3}{6}=\frac12.
\]

\textbf{Проверка.}
Из шести граней подходят ровно три. Это половина всех граней, поэтому вероятность $\frac12$ правдоподобна.

\outcome искомая вероятность равна $\frac12$.

\textbf{Задача на закрепление.}
Игральный кубик бросают один раз. Найдите вероятность того, что выпавшее число больше $3$ и при этом отлично от $6$.

\newpage
\phantomsection\label{task:5}
\section*{Задание 5}

\textbf{Текст задания.}
Монету бросают два раза. Найдите вероятность того, что орёл выпадет ровно один раз. Для решения можно выписать все возможные результаты опыта.

\textbf{Ваше решение.}
Записано: «не помню». Перечень исходов и вычисление вероятности отсутствуют.

\textbf{Ваш ответ.}
Не указан.

\errlabel{решение не выполнено.}

\textbf{Где именно возникает ошибка.}
Решение останавливается до построения модели двух последовательных бросков. Для такой задачи важно считать результатом опыта не отдельную сторону монеты, а упорядоченную пару результатов первого и второго бросков.

Последний корректный момент --- дано, что монету бросают два раза. Первый необходимый, но пропущенный шаг --- выписать все результаты этих двух бросков.

\textbf{Почему это важно.}
Обозначим орла буквой $O$, решку буквой $P$. За два броска возможны четыре равновероятных последовательности:
\[
OO,\ OP,\ PO,\ PP.
\]
Условие «ровно один раз» означает: в последовательности должен быть один орёл и одна решка. Поэтому подходят только
\[
OP,\ PO.
\]
Последовательность $OO$ не подходит, потому что орёл выпал два раза; последовательность $PP$ не подходит, потому что орёл не выпал ни разу.

\keyidea при нескольких последовательных испытаниях исход должен описывать весь опыт целиком. Слово «ровно» задаёт точное количество наступлений события, а не «хотя бы одно».

\textbf{Исправление.}
Всего равновероятных исходов $4$, благоприятных $2$. Поэтому
\[
P=\frac{2}{4}=\frac12.
\]

\textbf{Полное правильное решение.}

1. Выпишем все результаты двух бросков:
\[
OO,\ OP,\ PO,\ PP.
\]

2. Выберем те, где орёл встречается ровно один раз:
\[
OP,\ PO.
\]

3. Благоприятных исходов $2$, всего исходов $4$. Значит,
\[
P=\frac24=\frac12.
\]

\textbf{Проверка.}
Два благоприятных исхода из четырёх составляют половину пространства исходов. Кроме того, события «ровно один орёл» и «результаты двух бросков различаются» здесь совпадают, а различными результаты бывают в двух случаях из четырёх.

\outcome вероятность того, что орёл выпадет ровно один раз, равна $\frac12$.

\textbf{Задача на закрепление.}
Монету бросают два раза. Найдите вероятность того, что результаты двух бросков будут различными.

\newpage
\phantomsection\label{task:6}
\section*{Задание 6}

\textbf{Текст задания.}
В непрозрачном мешке находятся $5$ белых, $3$ чёрных и $2$ серых одинаковых шара. Из мешка случайным образом выбирают один шар. Найдите вероятность того, что выбранный шар окажется не чёрным.

\textbf{Ваше решение.}
Записано: «не помню». Подсчёт общего числа шаров и подходящих шаров отсутствует.

\textbf{Ваш ответ.}
Не указан.

\errlabel{решение не выполнено.}

\textbf{Где именно возникает ошибка.}
Решение не доходит до первого вычислительного шага. Нужно определить общее количество равновозможных шаров и число шаров, удовлетворяющих условию «не чёрный».

Последний корректный момент --- данные о количестве шаров каждого цвета. Первый пропущенный шаг --- сложить количества шаров всех цветов.

\textbf{Почему это нужно сделать.}
Всего в мешке
\[
5+3+2=10
\]
шаров. Событие «шар не чёрный» включает два цвета: белый и серый. Таких шаров
\[
5+2=7.
\]
Следовательно, из $10$ одинаково возможных шаров подходят $7$.

Эту же задачу можно проверить через противоположное событие. Вероятность выбрать чёрный шар равна $\frac3{10}$, поэтому вероятность не выбрать чёрный шар должна составлять оставшуюся часть единицы.

\keyidea слово «не» удобно понимать как противоположное событие. Можно либо напрямую сосчитать все подходящие исходы, либо вычесть вероятность запрещённого события из $1$.

\textbf{Исправление.}
Прямой подсчёт даёт
\[
P=\frac{5+2}{5+3+2}=\frac7{10}.
\]

\textbf{Полное правильное решение.}

1. Найдём общее число шаров:
\[
5+3+2=10.
\]

2. Не чёрными являются белые и серые шары. Их количество:
\[
5+2=7.
\]

3. Вероятность выбрать не чёрный шар:
\[
P=\frac7{10}.
\]

\textbf{Проверка.}
Вероятность выбрать чёрный шар равна
\[
\frac3{10}.
\]
Тогда вероятность противоположного события:
\[
1-\frac3{10}=\frac{10}{10}-\frac3{10}=\frac7{10},
\]
что совпадает с прямым подсчётом.

\outcome искомая вероятность равна $\frac7{10}$.

\textbf{Задача на закрепление.}
В мешке находятся $4$ красных, $5$ синих и $3$ зелёных одинаковых шара. Случайно выбирают один шар. Найдите вероятность того, что выбранный шар окажется не синим.

\newpage
\phantomsection\label{task:8}
\section*{Задание 8}

\textbf{Текст задания.}
По итогам контрольной работы получена таблица:
\[
\begin{array}{c|cccc}
\text{Оценка} & 2 & 3 & 4 & 5 \\
\hline
\text{Количество учеников} & 1 & 4 & 7 & 3
\end{array}
\]
Найдите: 1) среднюю оценку; 2) медиану; 3) моду.

\textbf{Ваше решение.}
Промежуточные вычисления не показаны. Записано:
\[
1)\ 4,\qquad 2)\ 3{,}5,\qquad 3)\ \text{---}.
\]

\textbf{Ваш ответ.}
$1)\ 4;\quad 2)\ 3{,}5;\quad 3)$ не указан.

\errlabel{неверно найдены среднее и медиана; мода не указана.}

\textbf{Где именно возникает ошибка.}
Поскольку вычисления не записаны, нельзя надёжно определить, какой приём использовался. Первый достоверно неверный фрагмент --- ответ $4$ к пункту о среднем значении. Следующий неверный фрагмент --- $3{,}5$ как медиана. В третьем пункте обязательный ответ отсутствует.

Сохранить можно только корректно переписанные исходные данные таблицы. Дальше каждую статистическую характеристику нужно находить по своему правилу.

\textbf{Почему это неверно.}

\emph{Среднее арифметическое.} В таблице каждая оценка встречается разное число раз. Поэтому нельзя работать только с четырьмя числами $2,3,4,5$, игнорируя их частоты. Нужно учесть, что одна двойка даёт вклад $2\cdot1$, четыре тройки --- $3\cdot4$, семь четвёрок --- $4\cdot7$, а три пятёрки --- $5\cdot3$.

Всего учеников:
\[
1+4+7+3=15.
\]
Сумма всех оценок с учётом повторений:
\[
2\cdot1+3\cdot4+4\cdot7+5\cdot3=2+12+28+15=57.
\]
Поэтому среднее равно $\frac{57}{15}=3{,}8$, а не $4$.

\emph{Медиана.} Наблюдений $15$, то есть их нечётное число. Значит, после упорядочивания медиана --- единственный центральный элемент, стоящий на восьмом месте. Здесь первое место занимает оценка $2$; места со второго по пятое занимают оценки $3$; места с шестого по двенадцатое --- оценки $4$. Восьмой элемент равен $4$. Поэтому медиана равна $4$, а не $3{,}5$.

\emph{Мода.} Мода --- значение, встречающееся чаще других. Наибольшая частота в таблице равна $7$, и она соответствует оценке $4$. Значит, мода равна $4$.

\keyidea одна таблица частот содержит всю выборку в сжатом виде, но для среднего, медианы и моды используются разные действия: среднее учитывает каждое значение столько раз, сколько оно встречается; медиана определяется положением в упорядоченном ряду; мода определяется максимальной частотой.

\textbf{Исправление от последнего правильного шага.}
Сохраняем таблицу и отдельно вычисляем три характеристики.

Количество учеников:
\[
15.
\]
Сумма оценок:
\[
57.
\]
Среднее:
\[
\frac{57}{15}=\frac{19}{5}=3{,}8.
\]

Для медианы при $15$ наблюдениях берём восьмой элемент. Восьмой элемент попадает в группу оценок $4$, значит медиана равна $4$.

Максимальная частота равна $7$, и она соответствует оценке $4$, значит мода равна $4$.

\textbf{Полное правильное решение.}

1. Найдём объём выборки:
\[
1+4+7+3=15.
\]

2. Найдём сумму всех оценок с учётом частот:
\[
2\cdot1+3\cdot4+4\cdot7+5\cdot3=57.
\]

3. Средняя оценка:
\[
\frac{57}{15}=\frac{19}{5}=3{,}8.
\]

4. Для медианы определим центральное место. При $15$ результатах центральным является восьмой результат. Распределение мест такое:
\[
\begin{array}{c|c}
\text{Оценка} & \text{Места в упорядоченном ряду} \\
\hline
2 & 1 \\
3 & 2\text{--}5 \\
4 & 6\text{--}12 \\
5 & 13\text{--}15
\end{array}
\]
Восьмое место относится к оценке $4$, поэтому медиана равна $4$.

5. Для моды сравним частоты $1,4,7,3$. Наибольшая частота --- $7$, она соответствует оценке $4$. Следовательно, мода равна $4$.

\textbf{Проверка.}
Среднее $3{,}8$ находится между минимальной оценкой $2$ и максимальной $5$ и близко к $4$, что естественно: именно четвёрок больше всего. Медиана $4$ согласуется с расположением восьмого результата внутри семи четвёрок. Мода $4$ согласуется с максимальной частотой $7$.

\outcome средняя оценка равна $3{,}8$, медиана равна $4$, мода равна $4$.

\textbf{Задача на закрепление.}
По результатам контрольной работы получена таблица:
\[
\begin{array}{c|cccc}
\text{Оценка} & 2 & 3 & 4 & 5 \\
\hline
\text{Количество учеников} & 2 & 3 & 6 & 1
\end{array}
\]
Найдите среднюю оценку, медиану и моду.

\vfill
\begin{center}
\small Лёвин Артём Александрович эксклюзивно для Екатерины
\end{center}

\end{document}
